\section{Cohomología y denominadores}
\label{sec:cohomologia}

\fuente{A, §8, pp.~35--41, especialmente los Lemas~8.1,
8.3, 8.5--8.6 y las entradas de las Proposiciones~8.2 y 8.4;
B, §5, pp.~52--64, en el ejemplar de 73 páginas.
Se desarrollan los mecanismos de longitudes, congruencias,
selección de torcidos y elección de ventana. Las cotas
aritméticas que alimentan esos mecanismos se enuncian con
su alcance; la concordancia con Lean se delimita al final.}

La Nota~\ref{sec:proyecciones} mostró cómo un proyector racional
puede introducir denominadores. Las Notas~\ref{sec:periodos}
y \ref{sec:geometria} incorporaron sus costes a la comparación
de pendientes. Estudiaremos ahora la cantidad cohomológica
que acota una parte de esos costes y la selección de un
intervalo en el que todos ellos se controlan simultáneamente.

Distinguimos dos cuerpos de números: $E$ es el cuerpo
aritmético de la construcción y $K$ es el cuerpo de
coeficientes de los periodos. Escribimos $v\mid\ell$ para
un lugar de $E$ y $w\mid\ell$ para uno de $K$.
El entero $e\ge1$ será el índice de un torcido de Tate;
$t$ será su corte superior. Reservamos $a$ para la altura
de la capa diádica y $\nu$ para un nivel de una torre
$\Z_\ell$. Así, $m$ puede seguir designando el número de
levantamientos theta de la Nota~\ref{sec:periodos}.

\subsection{Longitud, orden y cambio del anillo de coeficientes}

Sea $\mathcal O$ un anillo de valoración discreta, con
uniformizante $\varpi$ y cuerpo residual $\kappa$.
La longitud $\length_{\mathcal O}M$ cuenta los factores
simples de una serie de composición de $M$; cada factor
es isomorfo a $\kappa$. Puede valer infinito.
En una sucesión exacta corta las longitudes se suman,
también como cantidades extendidas; véase \cite{stacks-longitud}.

Cuando $\mathcal O$ es el anillo de enteros de una extensión
finita de $\Q_\ell$, ponemos
\[
 e_{\mathcal O}=\ord_{\varpi}(\ell),\qquad
 |\kappa|=\ell^{f_{\mathcal O}}.
\]
Para un módulo de longitud finita, su serie de composición da
\begin{equation}
 \log|M|=f_{\mathcal O}\,\length_{\mathcal O}M\,\log\ell.
 \label{eq:cohom-longitud-orden}
\end{equation}
En particular, sobre $\Z_\ell$ la longitud es
$\log|M|/\log\ell$.

Si $A$ es un $\Z_\ell$-módulo finito, la descomposición
$A\simeq\bigoplus_i\Z_\ell/\ell^{a_i}\Z_\ell$ muestra que
\begin{equation}
 \length_{\mathcal O}(A\otimes_{\Z_\ell}\mathcal O)
 =e_{\mathcal O}\sum_i a_i
 =e_{\mathcal O}\,\length_{\Z_\ell}A.
 \label{eq:cohom-base-cambio}
\end{equation}
En efecto, $\ell^{a_i}\mathcal O=\varpi^{e_{\mathcal O}a_i}
\mathcal O$ hasta una unidad. Esta cuenta explica la
división por el índice de ramificación en las longitudes
normalizadas de A y B. El grado residual interviene,
en cambio, al pasar de longitud a logaritmo del orden.

\begin{lema}[Exponente y torsión de orden \(\ell\)]
\label{lem:cohom-exponente}
Sea $M$ un $\Z_\ell$-módulo, sean $r,g\in\N$, y supongamos
que $\ell^rM=0$ y que
$\dim_{\mathbb F_\ell}M[\ell]\le g$.
Entonces $M$ es finito y
\[
 \length_{\Z_\ell}M\le rg.
\]
\end{lema}

\begin{proof}
El caso $r=0$ da $M=0$. Para $1\le i\le r$,
la multiplicación por $\ell^{i-1}$ induce una inyección
\[
 M[\ell^i]/M[\ell^{i-1}]\longrightarrow M[\ell].
\]
Su núcleo es cero por la definición del denominador.
Cada cociente es un espacio vectorial sobre $\mathbb F_\ell$
de dimensión a lo sumo $g$. La filtración
$0=M[\ell^0]\subseteq\cdots\subseteq M[\ell^r]=M$
tiene $r$ cocientes finitos, con orden a lo sumo $\ell^g$.
La multiplicación de sus órdenes prueba finitud y la cota.
\end{proof}

Así, una cota del exponente y una cota de la torsión
de orden $\ell$ prueban conjuntamente la finitud.
Este orden lógico se conserva al utilizar las entradas
aritméticas de B, p.~60.

\subsection{Ideales de congruencia y coste de una proyección}

Ampliamos el lema de la Nota~\ref{sec:proyecciones} a varios
sistemas de autovalores. Sea $F$ el cuerpo de fracciones
de un anillo de valoración discreta $\mathcal O$.
Sea $X$ un conjunto finito, y consideremos
\[
 V=\bigoplus_{\pi\in X}V_\pi,\qquad
 m_\pi=\dim_FV_\pi.
\]
Una $\mathcal O$-subálgebra $\mathbb T\subseteq\mathcal O^X$
actúa diagonalmente: $u\in\mathbb T$ actúa por $u_\pi$
en $V_\pi$. Suponemos que contiene los escalares,
separa los elementos de $X$ y preserva un retículo pleno $M$.
Separar significa que para $\pi\ne\pi'$ existe $u\in\mathbb T$
con $u_\pi\ne u_{\pi'}$.

Para $Z\subseteq X$, sean $V_Z=\bigoplus_{\pi\in Z}V_\pi$
y $e_Z:V\to V_Z$ el proyector. Ponemos
\begin{equation}
 \begin{gathered}
 L=e_ZM,\qquad N=M\cap V_Z,\\
 \mathfrak J=\{u\in\mathbb T:u_{\pi'}=0
       \text{ para }\pi'\notin Z\},\qquad
 \mathfrak c_\pi=\{u_\pi:u\in\mathfrak J\}\quad(\pi\in Z).
 \end{gathered}
 \label{eq:cohom-ideales}
\end{equation}
Cada $\mathfrak c_\pi$ es un ideal no nulo de $\mathcal O$.
Para verlo, por cada $\pi'\notin Z$ elegimos un elemento
que separa $\pi$ de $\pi'$ y le restamos el escalar
dado por su coordenada $\pi'$. El producto de estos
elementos pertenece a $\mathfrak J$ y su coordenada $\pi$
es no nula. Si $Z=X$, sirve el elemento $1$.

\begin{proposicion}[Cota por ideales de congruencia]
\label{prop:cohom-congruencias}
Con estos datos,
\begin{equation}
 \length_{\mathcal O}(e_ZM/(M\cap V_Z))
 \le\sum_{\pi\in Z}m_\pi\,
       \length_{\mathcal O}(\mathcal O/\mathfrak c_\pi).
 \label{eq:cohom-cota-congruencias}
\end{equation}
\end{proposicion}

\begin{proof}
Si $Z$ es vacío, ambos miembros son cero.
En otro caso ordenamos $Z=\{\pi_1,\ldots,\pi_s\}$ y ponemos
\[
 W_i=\bigoplus_{k\le i}V_{\pi_k},\qquad
 L_i=L\cap W_i,\qquad N_i=N\cap W_i.
\]
Los módulos $L$ y $N$ son retículos plenos de $V_Z$,
con $N\subseteq L$. Cada $L_i$ es pleno en $W_i$,
y $L_{i-1}$ es saturado en $L_i$ porque se obtiene
intersecando con un subespacio. Por tanto
$U_i=L_i/L_{i-1}$ es libre de rango $m_{\pi_i}$.

Elegimos $u_i\in\mathfrak J$ cuya coordenada
$a_i=(u_i)_{\pi_i}$ genera $\mathfrak c_{\pi_i}$.
Como $u_ie_Z=u_i$ y $u_iM\subseteq M$, se tiene
$u_iL_i\subseteq N_i$. En $U_i$, el operador $u_i$
actúa por $a_i$: las otras coordenadas de $W_i$
desaparecen al tomar el cociente. En consecuencia,
\[
 a_iU_i\subseteq (N_i+L_{i-1})/L_{i-1}.
\]
El cociente $L_i/(N_i+L_{i-1})$ tiene longitud a lo sumo
$m_{\pi_i}\ord_\varpi(a_i)$.

La filtración $(L_i+N)/N$ de $L/N$ tiene precisamente
esos cocientes graduados, pues
$L_i\cap(L_{i-1}+N)=L_{i-1}+N_i$.
Sumar sus longitudes prueba
\eqref{eq:cohom-cota-congruencias}.
\end{proof}

La prueba permite elegir un operador distinto en cada
graduado. Es válida aunque $u_i$ actúe de manera no nula
en otros espacios de $Z$.
B, p.~53, utiliza esta desigualdad para convertir cotas
de ideales de congruencia en cotas de denominadores.
La construcción de esos ideales a partir de las
representaciones automorfas constituye la entrada aritmética.

\subsection{Extensiones, cociclos y torcidos negativos}

Sea $G$ un grupo profinito y $A$ un módulo discreto
con acción continua de $G$. Un cociclo continuo es
una aplicación $z:G\to A$ que satisface
\begin{equation}
 z(gh)=z(g)+g\,z(h).
 \label{eq:cohom-cociclo}
\end{equation}
Los cociclos de la forma $z(g)=g\,a-a$ son cobordes;
su cociente define $H^1(G,A)$. El grupo
$H^0(G,A)=A^G$ está formado por los elementos invariantes.
Las clases de $H^1$ describen el defecto de compatibilidad
de una elección de secciones con la acción de $G$.

El origen del cociente de caracteres se ve en una matriz.
Sobre un anillo de coeficientes, una extensión de dos
módulos libres de rango uno puede escribirse, después
de elegir una sección del módulo cociente, como
\[
 \rho(g)=
 \begin{pmatrix}\alpha(g)&b(g)\\0&\beta(g)\end{pmatrix}.
\]
La identidad $\rho(gh)=\rho(g)\rho(h)$ implica que
$z(g)=b(g)\beta(g)^{-1}$ satisface
\[
 z(gh)=z(g)+\alpha(g)\beta(g)^{-1}z(h).
\]
Cambiar la sección por un múltiplo $a$ del primer vector
cambia $z(g)$ en
$z(g)+(\alpha(g)\beta(g)^{-1}-1)a$.
La acción en el módulo de extensiones es, por tanto,
la del carácter $\alpha\beta^{-1}$.

En B, §5.2, pp.~54--55, las $h$ líneas añadidas por la torre
tienen caracteres, para $1\le c\le h$,
\[
 \xi_c=\xi_1\chi_\ell^{-(c-1)}.
\]
Para $c>d$, la razón $\xi_c\xi_d^{-1}$ es
$\chi_\ell^{-(c-d)}$. Así aparecen los torcidos negativos
de índices $1,\ldots,h-1$. La selección de bloques de
extensiones y sus condiciones locales forman parte de la
comparación automorfa; el cálculo matricial identifica
sus coeficientes.

Escribiremos
\begin{equation}
 A_{\ell,e}=\Q_\ell/\Z_\ell(-e),\qquad
 g\cdot a=\chi_\ell(g)^{-e}a,\qquad e\ge1.
 \label{eq:cohom-torcido}
\end{equation}
El grupo aditivo subyacente es $\Q_\ell/\Z_\ell$;
el torcido especifica su acción de Galois.

\subsection{El núcleo local y la corrección de dualidad}

Sea $S_\ell$ el conjunto de lugares finitos de $E$
por encima de $\ell$. El grupo $G_{E,S_\ell}$ corresponde
a la extensión máxima no ramificada fuera de esos lugares
y los arquimedianos. Definimos
\begin{equation}
 \operatorname{Sha}_{\ell,e}(E)=
 \ker\left(
 H^1(G_{E,S_\ell},A_{\ell,e})
 \longrightarrow\prod_{v\mid\ell}H^1(E_v,A_{\ell,e})
 \right).
 \label{eq:cohom-sha}
\end{equation}
Escribimos $s_{\ell,e}(E)=
\length_{\Z_\ell}\operatorname{Sha}_{\ell,e}(E)$,
con valores en $\N\cup\{\infty\}$.
Todas las cohomologías con estos coeficientes discretos
son continuas. Para coeficientes $\Z_\ell(e+1)$ se usa
la interpretación continua correspondiente al límite
de los coeficientes finitos.

La localización de \eqref{eq:cohom-sha} contiene los
lugares finitos sobre $\ell$. Para $v\nmid\ell$,
una clase global es no ramificada. Si $q_v$ es el tamaño
del cuerpo residual, Frobenius geométrico actúa en
$A_{\ell,e}$ por $q_v^e$. La descripción de la primera
cohomología no ramificada como
$A_{\ell,e}/(q_v^e-1)A_{\ell,e}$ muestra que es cero:
$q_v^e-1\ne0$ y la multiplicación por cualquier entero
no nulo es sobreyectiva en $\Q_\ell/\Z_\ell$.
Cuando $E$ es CM, sus completaciones arquimedianas
son complejas y su primera cohomología también es cero.
Esta es la observación de B, Remark~5.2, p.~55.

Recibimos de la dualidad de Poitou--Tate la forma adaptada
a estos coeficientes, incluidos los pasos a límites
y las condiciones en los lugares reales.
Como referencia del teorema para módulos finitos,
véase \cite{milne-dualidad}, I, §4, Teorema~4.10, p.~57.
En la situación CM de A/B, el fragmento relevante es
\begin{equation}
 \begin{gathered}
 0\longrightarrow H^0(G_{E,S_\ell},A_{\ell,e})
 \longrightarrow\prod_{v\mid\ell}H^0(E_v,A_{\ell,e})\\
 \longrightarrow
 H^2(G_{E,S_\ell},\Z_\ell(e+1))^\vee
 \longrightarrow\operatorname{Sha}_{\ell,e}(E)
 \longrightarrow0.
 \end{gathered}
 \label{eq:cohom-dualidad}
\end{equation}
Aquí el superíndice \({}^\vee\) indica el dual
de Pontryagin, que para un grupo finito conserva su
orden y su longitud. La teoría de dualidad, y no sólo
el cálculo de cociclos, proporciona la exactitud
de este fragmento.

La sobreyección desde el dual de $H^2$ hacia Sha
ya muestra que $H^2$ finito implica Sha finito.
Así se obtiene la finitud del núcleo sin presuponerla
para escribir una identidad de longitudes finitas.

Supongamos ahora que los términos tienen longitud finita.
Separar \eqref{eq:cohom-dualidad} en dos sucesiones
exactas cortas da
\begin{equation}
 \begin{aligned}
 &\length_{\Z_\ell}H^2(G_{E,S_\ell},\Z_\ell(e+1))
   +\length_{\Z_\ell}H^0(G_{E,S_\ell},A_{\ell,e})\\
 &\hspace{1em}=
 s_{\ell,e}(E)+
 \sum_{v\mid\ell}\length_{\Z_\ell}H^0(E_v,A_{\ell,e}).
 \end{aligned}
 \label{eq:cohom-longitudes-dualidad}
\end{equation}
La primera flecha es inyectiva, por lo que la suma
de longitudes locales es al menos la longitud global.
Resulta
\begin{equation}
 s_{\ell,e}(E)\le
 \length_{\Z_\ell}H^2(G_{E,S_\ell},\Z_\ell(e+1)).
 \label{eq:cohom-sha-le-hdos}
\end{equation}
Esta es la dirección que permite utilizar una cota
superior de $H^2$ para estimar el núcleo local.
La comparación también se conserva como desigualdad
de longitudes extendidas antes de probar finitud.
En la dirección inversa hacen falta las correcciones
$H^0$. A, (8.8), p.~40, aporta una cota específica de
ellas al construir la entrada de la capa fija.

\subsection{Un presupuesto simultáneo de longitudes}

Para un entero $t\ge1$ definimos
\begin{equation}
 \mathcal V(E,t)=
 \sup_{\mathcal P\text{ finito}}
 \sum_{\ell\in\mathcal P}(\log\ell)
       \sum_{1\le e<t}s_{\ell,e}(E),
 \label{eq:cohom-presupuesto}
\end{equation}
donde $\mathcal P$ recorre conjuntos de primos.
Si alguna longitud requerida es infinita, el valor es infinito;
si $t=1$, la suma es vacía y vale cero.
La notación $\mathcal V(E,t;B)$ de A, (8.2), equivale a
afirmar finitud de todas las longitudes y
$\mathcal V(E,t)\le B$.
Escribimos las sumas primero sobre conjuntos finitos para
no presuponer su convergencia.

\begin{lema}[Sumas parciales y corte]
\label{lem:cohom-presupuesto}
Si todas las longitudes requeridas son finitas y todas las
sumas parciales de \eqref{eq:cohom-presupuesto} son a lo
sumo $B<\infty$, entonces $\mathcal V(E,t)\le B$.
Además, $1\le t'\le t$ implica
$\mathcal V(E,t')\le\mathcal V(E,t)$.
Una partición de los primos en un número finito de clases
con presupuestos $B_1,\ldots,B_u$ da el presupuesto total
$B_1+\cdots+B_u$.
Hay a lo sumo $\lfloor B/\log2\rfloor$ pares $(\ell,e)$
con $1\le e<t$ y $s_{\ell,e}(E)>0$.
\end{lema}

\begin{proof}
La primera afirmación es la definición del supremo.
Para la segunda, cada suma con corte $t'$ es una subsuma
con sumandos no negativos de la correspondiente con corte $t$.
Para la tercera, se intersecta cualquier conjunto finito de
primos con las clases y se suman sus cotas.
Finalmente, cada longitud positiva es un entero al menos uno,
así que cada par positivo aporta al menos $\log2$.
Cualquier colección finita de esos pares queda contenida en
una suma parcial; su cardinal es, por tanto, a lo sumo
$B/\log2$. Esto acota también la colección completa.
\end{proof}

La finitud de cada grupo, por sí sola, no da ese presupuesto.
La familia artificial $s_{\ell,1}=1$ para todos los primos
tiene cada término finito y sumas parciales arbitrariamente
grandes. En el cotejo formal, esta distinción reaparecerá
como una hipótesis de sumabilidad explícita.

Para $E$ CM abeliano, la conjugación compleja induce una
involución $c$ en cohomología, mediante un levantamiento
que normaliza $G_E$ y su acción en los coeficientes.
No se la considera un elemento de $G_E$ que fijara $E$.
Cuando $\ell$ es impar, B, p.~58, usa las etiquetas
\[
 \operatorname{Sha}_{\ell,e}^{+}
   =\ker(c-(-1)^e),\qquad
 \operatorname{Sha}_{\ell,e}^{-}
   =\ker(c-(-1)^{e+1}).
\]
Los operadores
\[
 P_+=\frac{1+(-1)^ec}{2},\qquad
 P_-=\frac{1-(-1)^ec}{2}
\]
son idempotentes complementarios: se comprueba expandiendo
sus productos, con $c^2=1$ y $2$ invertible.
Sus imágenes dan una suma directa y
$s_{\ell,e}=s_{\ell,e}^{+}+s_{\ell,e}^{-}$.
Para $e$ impar, el signo $+$ designa aquí el autoespacio
$c=-1$. El torcido forma parte de la convención.

En $\ell=2$ no se dispone de esos proyectores integrales.
Por ejemplo, en $\Z/4\Z$ con $c=-1$, el autoespacio
$c=+1$ es $2\Z/4\Z$ y el autoespacio $c=-1$ es todo el
módulo: no constituyen una suma directa.
Por eso la contribución diádica se trata separadamente.
Si $B_2$ acota $\sum_{1\le e<t}s_{2,e}(E)$ y
$B_\pm$ acotan las sumas impares correspondientes,
ya ponderadas por $\log\ell$, resulta
\begin{equation}
 \mathcal V(E,t)\le B_2\log2+B_-+B_+.
 \label{eq:cohom-particion}
\end{equation}
El factor $\log2$ restaura el peso que se omitió al
definir $B_2$.

\subsection{Del índice local a la pérdida de pendiente}

Volvamos a la etapa theta $j$ de las notas anteriores.
Escribimos $M^{\rm int}_{j,w}\subseteq M^{\rm pr}_{j,w}$
para el retículo de intersección y el proyectado, de igual
rango $d_j>0$ sobre $\mathcal O_w$.
Con $e_w$ índice de ramificación de $K_w/\Q_\ell$,
B, pp.~55--56, define
\begin{equation}
 \begin{aligned}
 \eta_{j,w}
   &=\frac{1}{e_w}\length_{\mathcal O_w}
       (M^{\rm pr}_{j,w}/M^{\rm int}_{j,w}),\\
 \eta_j
   &=\frac{1}{[K:\Q]}
       \sum_w e_w\eta_{j,w}\log Nw.
 \end{aligned}
 \label{eq:cohom-indices}
\end{equation}
La construcción proporciona un conjunto finito de lugares
donde cambia la elección local.
La primera división normaliza la longitud por ramificación;
el segundo $e_w$ restituye la longitud integral al calcular
el índice global. Aquí no aparece un factor dos.

Si $M^{\rm int}_j\subseteq M^{\rm pr}_j$ son los retículos
globales correspondientes y conservan las mismas métricas,
la cuenta de cardinalidades por lugares y la normalización
de pendientes de la Nota~\ref{sec:marcos} dan
\begin{equation}
 \frac{\log\#(M^{\rm pr}_j/M^{\rm int}_j)}{[K:\Q]}
 =\eta_j,\qquad
 \mu(M^{\rm int}_j)
 =\mu(M^{\rm pr}_j)-\frac{2\eta_j}{d_j}.
 \label{eq:cohom-coste-pendiente}
\end{equation}
Así se distingue la suma de defectos $\eta_j$ del coste
de pendiente $2\eta_j/d_j$ y del parámetro espectral
del marco de la Nota~\ref{sec:geometria}.
Si se divide posteriormente por un rango retenido menor,
debe conservarse el cociente entre ambos rangos.

También se comprueba directamente la desaparición del
cuerpo de coeficientes en una cota uniforme local.
Si $\eta_{j,w}\le d_jc_\ell$ para todos $w\mid\ell$,
con $c_\ell\ge0$, entonces
\[
 \frac{\eta_j}{d_j}
 \le\frac{1}{[K:\Q]}
     \sum_\ell c_\ell\log\ell\sum_{w\mid\ell}e_wf_w
 =\sum_\ell c_\ell\log\ell,
\]
pues $\log Nw=f_w\log\ell$ y
$\sum_{w\mid\ell}e_wf_w=[K:\Q]$.
Un subconjunto de lugares se completa usando no negatividad.

\paragraph{La entrada aritmética.}
La comparación especializada entre congruencias y extensiones
es lo que permite aplicar esta cuenta.
B, (5.1), p.~56, afirma, con sus parámetros y suponiendo
$\mathcal V(E,h)<\infty$,
\begin{equation}
 \eta_j\le d_j\left[
 C_2b x^{2\varepsilon}
 \left(fm+\frac{|T_2|k(\log x)^2}{n_{\rm sel}^{\,2}}\right)
 +\mathcal V(E,h)\right].
 \label{eq:cohom-entrada-indice}
\end{equation}
Aquí $h$ es el entero de la torre theta, $m$ su número de
levantamientos, $f=[E^+:\Q]$ y $n_{\rm sel}$ designa el
$n=|\mathcal H|$ de B, pp.~54--55.
No se lo identifica con $[E:\Q]$.
Los parámetros $b,k,T_2$ conservan el significado de la
construcción de las Notas~\ref{sec:periodos} y
\ref{sec:geometria}; $C_2$ se fija antes de hacerlos variar.
El término de longitudes sí queda definido en esta nota.
Los otros costes y el puente aritmético siguen siendo
entradas de la construcción.

La representación recibida en B, pp.~53--55, tiene un
sumando $\rho_\lambda$ y las líneas
$\xi_c=\xi_1\chi_\ell^{-(c-1)}$.
Los grupos de extensiones entre líneas llevan los torcidos
negativos ya calculados. Para el sumando $\rho_\lambda$
se utiliza, en cambio, un grupo de Selmer con condición
finita de Bloch--Kato sobre $\ell$, condición no ramificada
fuera de $\ell$ y de los lugares adicionales permitidos,
y sin restricción en esos lugares adicionales.
No es el mismo núcleo de anulación local de
\eqref{eq:cohom-sha}.

Si $T_\pi$ es el retículo estable del sumando torcido,
escribimos $A_\pi=(T_\pi\otimes K_\lambda)/T_\pi$.
La elección especializada de caracteres proporciona en B
finitud y la longitud normalizada
\[
 s_\pi=
 \frac{\length_{\mathcal O_\lambda}
          \operatorname{Sel}(E,A_\pi)}{e_\lambda}
 \le\nu_\ell+\ell^{-2},\qquad
 \nu_\ell=\min_{g\in G_E}
      v_\ell\bigl(\chi_\ell(g)^{hr}-1\bigr),
\]
con $r=\dim\rho_\lambda$ y $v_\ell(\ell)=1$, además de
\[
 \sum_\ell\nu_\ell\log\ell
 \le2\log n_{\rm sel}+\log h+\log r.
\]
Estas cotas, con su finitud y sus condiciones locales, son
parte de la entrada.
El resto $\sum_\ell\ell^{-2}\log\ell$ sí se suma
elementalmente: lo domina
$\sum_{n\ge2}(\log n)/n^2<\infty$.
La finitud de un Selmer cualquiera no implicaría una cota
uniforme de esta forma.

\subsection{Los cuerpos variables y el orden de las elecciones}

Fijemos un carácter primitivo $\psi$, su cuerpo $F_\psi$,
su conductor $N_\psi$ y $o=[F_\psi:\Q]$.
Para primos auxiliares $q,d$ y altura $a\ge2$ se toma
\begin{equation}
 \begin{aligned}
 E_a&=F_\psi\,\Q(\zeta_{2^{a+2}})^+\,
          \Q(\sqrt{-q})\,\Q(\sqrt{-d}),\\
 N_a&=\operatorname{lcm}(N_\psi,2^{a+2})qd.
 \end{aligned}
 \label{eq:cohom-campos}
\end{equation}
El conductor $F_a$ de $E_a$ divide $N_a$.
Los factores pueden intersectarse; no se multiplican sus
grados como una igualdad. Sí tenemos
$[E_a:\Q]\le 4o\,2^a$ y
$[E_a:F_\psi]$ potencia de dos.
En consecuencia, todo primo impar que divide $[E_a:\Q]$
divide $o$, independientemente de $a,q,d$.

Para una escala $Y_0>1$, la admisibilidad de B, pp.~56--57,
incluye
\[
 q\equiv3\pmod8,\qquad d\equiv7\pmod{16},\qquad q\mid d+1,
\]
así como $p\mid d+1$ para todo primo impar
$p\le Y_0^{\varepsilon/5}$, y los intervalos de tamaño
\begin{equation}
 \begin{aligned}
 \tfrac12Y_0^{\varepsilon/5}
   &<\log q\le Y_0^{\varepsilon/5},\\
 \tfrac12Y_0^{\varepsilon/5}(\log Y_0)^2
   &<\log d\le Y_0^{\varepsilon/5}(\log Y_0)^2.
 \end{aligned}
 \label{eq:cohom-primos}
\end{equation}
La existencia simultánea de estos primos es una entrada
de la construcción; el argumento de esta nota trabaja con
cualquier pareja admisible.
Las cotas inferiores también se usan: para un umbral
suficientemente grande aseguran $q,d>2X^{401}+1$ cuando
$Y_0=X^{1/J}$ y $J$ ya está fijado.
En efecto, $X^{\varepsilon/(5J)}$ acaba dominando
cualquier múltiplo fijo de $\log X$.

Se fija además $0<\varepsilon<1/52$ y una función entera
$a(x)$ con
\begin{equation}
 x^{2+8\varepsilon}\le2^{a(x)}
       \le2x^{2+8\varepsilon}.
 \label{eq:cohom-altura}
\end{equation}
La meta es obtener $J\ge1$ y $X_0$ tales que para todo
$X\ge X_0$ y toda pareja admisible a escala $X^{1/J}$
exista $W$ con
\begin{equation}
 X^{1/J}\le W,\quad W^2\le X,\qquad
 \mathcal V(E_{a(x)},\lfloor x^{401}\rfloor)
 \le \exp\!\left(\frac{x^\varepsilon}{2}\right)
 \quad(W\le x\le W^2).
 \label{eq:cohom-meta}
\end{equation}
La misma pareja $q,d$ se conserva en todo el intervalo.
La altura puede cambiar con $x$; no se necesita que sea
monótona. Los parámetros uniformes se eligen antes de
$X,q,d$, mientras que la posición de la ventana puede
depender de ellos. Esta es la organización de A,
Teorema~8.7, p.~41, y B, Proposición~5.3, p.~57.

\subsection{La parte menos y los productos enteros}

El siguiente paso transforma muchas valuaciones locales en
una sola cantidad arquimediana. Es el Lema~8.1 de A, p.~36.

\begin{lema}[Producto entero]
\label{lem:cohom-producto}
Sea $M\ne0$ un entero y supongamos
$0\le b_\ell\le\ord_\ell(M)$ para los primos considerados.
Entonces toda suma parcial satisface
\[
 \sum_{\ell\in\mathcal P}b_\ell\log\ell\le\log|M|.
\]
Si $M$ es el producto de una colección finita, con
multiplicidades, de enteros algebraicos no nulos estable
bajo Galois, se obtiene la misma conclusión a partir
de $b_\ell\le\sum_\alpha v_\ell(\alpha)$, con
$v_\ell(\ell)=1$.
\end{lema}

\begin{proof}
La factorización prima da
$\log|M|=\sum_\ell\ord_\ell(M)\log\ell$, con sumandos
no negativos. En el segundo caso, el producto es racional
por invariancia de Galois, entero por integridad y no nulo
porque todos los factores lo son.
La multiplicatividad de la valoración da
$v_\ell(M)=\sum_\alpha v_\ell(\alpha)$.
\end{proof}

Si un primer producto $M_1$ permite la comparación salvo
en un primo $r_1$, y un segundo $M_2$ salvo en otro
$r_2\ne r_1$, usamos el primero fuera de $r_1$ y el segundo
en $r_1$. El presupuesto es a lo sumo $\log|M_1M_2|$.
Esta es una función concreta de los dos primos auxiliares.
La no anulación, la integridad y la estabilidad de los
factores deben establecerse antes de usar una cota de
sus valores absolutos complejos.

\paragraph{La entrada de valores especiales.}
Sea $E$ CM abeliano de grado $n_E$ y conductor $F$.
Para $e\ge1$ se consideran los caracteres $\theta$ de
$\operatorname{Gal}(E/\Q)$ con
$\theta(c)=(-1)^{e+1}$ y el carácter primitivo
$\chi_\theta$ de conductor $f_\theta\mid F$ que inducen.
B, p.~59, utiliza
\begin{equation}
 \begin{aligned}
 B_{e+1,\chi_\theta}
   &=f_\theta^{\,e}\sum_{u=1}^{f_\theta}
       \chi_\theta(u)B_{e+1}(u/f_\theta),\\
 \beta_{e,\theta}
   &=F\left(\prod_{\substack{p\le e+2\\p\ {\rm primo}}}p\right)
        B_{e+1,\chi_\theta}.
 \end{aligned}
 \label{eq:cohom-bernoulli}
\end{equation}
Aquí $B_k(T)$ es el polinomio de Bernoulli.
La fuente usa los resultados clásicos sobre sus
denominadores y sobre valores especiales para que
$\beta_{e,\theta}$ sea entero algebraico no nulo.
La paridad evita el cero trivial; el valor
$L(-e,\chi_\theta)=-B_{e+1,\chi_\theta}/(e+1)$ se relaciona
por la ecuación funcional con un valor de parte real
$e+1>1$. No se introduce GRH en este paso.

La eliminación de denominadores en
\eqref{eq:cohom-bernoulli} puede verse expandiendo
\[
 f_\theta^{\,e}B_{e+1}(u/f_\theta)
 =\sum_{r=0}^{e+1}\binom{e+1}{r}
       B_r u^{e+1-r}f_\theta^{\,r-1}.
\]
El único denominador de conductor es $f_\theta$ en
$r=0$, eliminado por $F$.
Los denominadores clásicos de los $B_r$ son eliminados
por el producto de primos hasta $e+2$.
La colección de factores es estable bajo Galois,
que permuta los caracteres y conserva su paridad.
Por tanto $M_e=\prod_\theta\beta_{e,\theta}$ es un
entero no nulo.

La afirmación aritmética sustantiva de B es, para
$\ell\nmid2n_E$,
\begin{equation}
 s_{\ell,e}^{-}(E)
 \le\sum_\theta v_\ell(\beta_{e,\theta}).
 \label{eq:cohom-menos-entrada}
\end{equation}
La construcción de anuladores y la comparación con
cohomología producen esa desigualdad; la cuenta de
Bernoulli no la demuestra.
Recibida la entrada, el
Lema~\ref{lem:cohom-producto} acota la suma sobre todos
esos primos por $\log|M_e|$.

La estimación arquimediana utilizada en B es
$|\beta_{e,\theta}|\le16F^{e+2}(e+1)!$.
Procede de la cota clásica de Fourier para los polinomios
de Bernoulli y del producto de primos hasta $e+2$.
Hay a lo sumo $n_E$ caracteres. Para $t\ge2$ y
$1\le e<t$,
\[
 \log|M_e|
 \le n_E\log\!\bigl(16F^{e+2}(e+1)!\bigr)
 \le n_E(2t+5)\log(tF).
\]
En la última cuenta se usan
$16\le(tF)^4$, $e+2\le t+1$ y $(e+1)!\le t^t$.
Como $(t-1)(2t+5)\le3t^2$, se obtiene
\begin{equation}
 \sum_{1\le e<t}
 \sum_{\substack{\ell\ {\rm impar}\\\ell\nmid n_E}}
   s_{\ell,e}^{-}(E)\log\ell
 \le3n_Et^2\log(tF).
 \label{eq:cohom-menos-buenos}
\end{equation}

\paragraph{Los primos que dividen el grado.}
Los proyectores de caracteres pueden tener denominadores
en $\ell\mid n_E$.
B, p.~60, utiliza entonces un anulador integral en un
anillo de grupo. La cuenta de adjugata explica cómo
se extrae un exponente: si la multiplicación por $z$
en un álgebra finita libre tiene determinante
$\Delta\ne0$, existe $y$ con $zy=\Delta$.
Si $z$ anula $bS$, entonces $\Delta b$ anula $S$.
En efecto, se aplica la identidad matriz por adjugata
a la coordenada del elemento unidad del álgebra.

La entrada específica proporciona un exponente $r$ y
una cota de torsión
\[
 \ell^r\operatorname{Sha}_{\ell,e}^{-}(E)=0,\quad
 r\log\ell\le C_0 n_E(e+1)\log(e\ell F),\quad
 \dim_{\mathbf F_\ell}
     \operatorname{Sha}_{\ell,e}^{-}(E)[\ell]
 \le8\ell n_E\log(\ell F).
\]
El Lema~\ref{lem:cohom-exponente} prueba primero la finitud
y después da, al sumar en $e<t$,
\[
 \sum_{1\le e<t}s_{\ell,e}^{-}(E)\log\ell
 \le8C_0\ell n_E^2t^2\log(t\ell F)^2.
\]
No bastaría solo el exponente ni solo la dimensión de
la torsión de orden $\ell$.

Para $E=E_a$ de \eqref{eq:cohom-campos}, los primos impares
que dividen $n_E$ pertenecen al conjunto fijo de divisores
de $o$. Así se recupera A, Proposición~8.2, pp.~36--37,
y B, (5.4), p.~60:
\begin{equation}
 B_-\le3n_Et^2\log(tF_a)
   +C\sum_{\substack{\ell\mid o\\\ell\ne2}}
       \ell n_E^2t^2\log(t\ell F_a)^2.
 \label{eq:cohom-menos-total}
\end{equation}
Los finitos primos de la segunda suma se pueden absorber
en una constante dependiente de $\psi$.
Esa constante ya no se describe como absoluta.

\subsection{La parte más: casos, deformaciones y multiplicidades}

El bloque de B, pp.~60--62, presenta la entrada agregada
\begin{equation}
 B_+\le(F_at)^C+C B_{\rm exc},
 \label{eq:cohom-mas-total}
\end{equation}
con $C$ fijo para la familia.
$B_{\rm exc}$ controla las contribuciones de un número
fijo de torres excepcionales.
La construcción aritmética de esta desigualdad usa
realizaciones automorfas y problemas de deformación
con condiciones locales precisas.
Separamos aquí las cuentas elementales que explican su
forma de esas entradas especializadas.

A, pp.~37--39, fija un conjunto finito de primos impares
$\mathcal B_\psi$, que contiene $3$ y depende solo de
$\psi$ y de sus auxiliares fijos.
Para un carácter $\theta$ se quita primero su componente
de orden potencia de $\ell$ y después su componente de
conductor potencia de $\ell$.
El carácter resultante $\theta_\ell^\circ$ tiene orden
y conductor primos a $\ell$.
La separación par/impar siguiente se refiere al
\emph{orden} de ese carácter.
Llamamos resonancia a $\ell-1\mid2e$.

\begin{center}
\small
\begin{tabularx}{\textwidth}{@{}>{\raggedright\arraybackslash}p{0.8cm}YY@{}}
\toprule
Caso & Condición relevante & Entrada utilizada\\
\midrule
I & $\ell\notin\mathcal B_\psi$, sin resonancia
  & Cota del total sobre primos y caracteres mediante productos.\\
Z & $\ell\notin\mathcal B_\psi$, resonancia, $\theta=1$
  & Anulación del carácter trivial.\\
II & $e\ge2$, orden de $\theta_\ell^\circ$ par;
       primo excepcional o resonancia
   & Cota polinómica en $\ell N_ae$.\\
III & $\ell\notin\mathcal B_\psi$, resonancia;
       $\theta\ne1$ de orden impar
    & Cota polinómica en $\ell N_ae$.\\
IV & $\ell\in\mathcal B_\psi$, $e\ge2$;
      orden de $\theta_\ell^\circ$ impar
   & Reducción a una torre de base fija.\\
V & $\ell\in\mathcal B_\psi$, $e=1$
  & Dos ramas según el conductor diádico.\\
\bottomrule
\end{tabularx}
\end{center}

La tabla conserva el reparto de la fuente, sin reemplazar
sus hipótesis locales.
Fuera de $\mathcal B_\psi$, en resonancia se tiene
$\ell\le2e+1<q,d$ y $\ell\nmid N_\psi$; en la familia,
esto permite usar $\theta_\ell^\circ=\theta$.
Si además $e=1$, se tendría $\ell=3$, ya incluido en
el conjunto excepcional.
Dentro de él, se separa $e=1$ y luego la paridad del orden
para $e\ge2$. De este modo no se pierde el caso V.

\begin{lema}[Conteo por núcleos]
\label{lem:cohom-nucleos}
Si $A$ es un grupo abeliano finito, entonces
\[
 \sum_{\substack{H\le A\\A/H\ \text{cíclico}}}
     \varphi([A:H])=|A|.
\]
\end{lema}

\begin{proof}
Se particionan los caracteres de $A$ según su núcleo.
La imagen de cada carácter es cíclica. Para un núcleo
$H$, los caracteres fieles del cociente cíclico $A/H$
son exactamente $\varphi([A:H])$.
La Nota~\ref{sec:caracteres} probó que hay $|A|$
caracteres en total; sumar la partición da la igualdad.
\end{proof}

La cota por núcleo de A, (8.5), p.~38, lleva ese factor
$\varphi([A:H])$. El lema controla el coste total de
agrupar los caracteres después de recibir la estimación
aritmética por núcleo.
En el caso I, la construcción de un producto entero no
nulo es lo que permite sumar primos sin cota superior.
Solo después, en los casos resonantes restantes,
$\ell\le2e+1<2t+1$ deja a lo sumo $2t$ primos, $t$
torcidos y $[E_a:\Q]\le N_a$ caracteres.
Una cota por término de la forma $(\ell N_ae)^C$, con
el peso $\log\ell$, se convierte así en otra potencia
fija de $N_at$.

\paragraph{Qué aporta la cohomología a una deformación.}
Sea $\bar\rho:G\to\operatorname{GL}_n(k)$ y
$k[\delta]/(\delta^2)$ el anillo de números duales.
Una deformación que reduce a $\bar\rho$ se escribe
$\rho_\delta(g)=(I+\delta c(g))\bar\rho(g)$.
Al multiplicar matrices se obtiene
\[
 c(gh)=c(g)+\bar\rho(g)c(h)\bar\rho(g)^{-1}.
\]
Es la ecuación de cociclo para la representación adjunta.
Conjugar por $I+\delta X$ cambia $c(g)$ por
$c(g)+X-\bar\rho(g)X\bar\rho(g)^{-1}$, un coborde.
Así aparece $H^1(G,\operatorname{ad}\bar\rho)$ como
espacio de deformaciones de primer orden.
Fijar determinante, polarización o condiciones locales
restringe ese espacio.
Este cálculo no demuestra las comparaciones entre
anillos de deformación y álgebras de Hecke, ni sus
estimaciones cuantitativas.

Dos cuentas adicionales aclaran las distintas entradas.
Para $r\ge31$, $e\ge1$ y $|z|=1$,
\[
 |zr^e+1+z^{-1}r^{-e}|
 \ge r^e-1-r^{-e}>3r^{e-1/2}.
\]
La segunda desigualdad se sigue de
$\sqrt r-r^{-1/2}-r^{-3/2}>3$.
Esa traza no puede coincidir con una traza cuyo valor
absoluto sea a lo sumo $3r^{e-1/2}$.
La realización automorfa, la cota de la segunda traza y
el producto integral que se construye con las diferencias
siguen siendo entradas.

Por otra parte, los cuatro caracteres
$\alpha\Psi\oplus\Psi\oplus\alpha\oplus1$ preservan la
forma alternante
$x_1y_4-x_4y_1+x_2y_3-x_3y_2$ con multiplicador
$\alpha\Psi$.
Para $\alpha=\chi_\ell$ y $\Psi=\theta\chi_\ell^e$,
los exponentes son $e+1,e,1,0$.
Son distintos cuando $e\ge2$, mientras que $e=1$
repite los dos exponentes centrales.
Esto explica por qué una entrada de peso regular no
resuelve por sí sola el caso V.
B, pp.~61--62, recurre a comparaciones adicionales de
Selmer y a dos torcidos concretos en el argumento de
anulación racional. Su alcance no se extiende a cualquier
grupo de Selmer.

\subsection{La entrada de capa fija y el papel de la homología}

A, Proposición~8.4, pp.~39--40, y B, p.~64, utilizan
la siguiente entrada, denominada \texttt{LemmaK} en el
código. Existe un exponente absoluto $C_K\ge1$ tal que,
para un cuerpo totalmente complejo $L$ que contiene
$\mu_\ell$ si $\ell$ es impar o $\mu_4$ si $\ell=2$,
y para $e\ge1$, el grupo
$H^2(G_{L,S_\ell},\Z_\ell(e+1))$ es finito y
\begin{equation}
 \begin{split}
 (\log\ell)\,W_\ell(L,e)
 &=
 \log\#H^2(G_{L,S_\ell},\Z_\ell(e+1))\\
 &\le
 \exp\!\left((2e[L:\Q])^{C_K}
                   \log(2|\operatorname{disc}L|)\right),
 \end{split}
 \label{eq:cohom-lema-k}
\end{equation}
donde $W_\ell(L,e)$ denota la longitud de ese $H^2$
sobre $\Z_\ell$.
Lo acotado es el \emph{logaritmo} de la cardinalidad.
La finitud forma parte de la entrada; no se deduce
simplemente asignando una expresión real a una longitud
que todavía pudiera ser infinita.

La ruta que A desarrolla para esta cota tiene varios
ingredientes. Una entrada de homología impar, llamada
\texttt{BOdd}, controla el rango de una imagen filtrada
especificada por operadores de bloques y estabilización.
No afirma que toda la homología impar tenga ese rango.
Las etapas \texttt{SClass}, \texttt{SDetect},
\texttt{SFew} y \texttt{SCount} construyen clases,
las detectan y comparan su número con caracteres
equivariantes de grupos de clases.
La torsión integral, los pasos al límite y el descenso
a cuerpos auxiliares completan la comparación con
cohomología.
Son entradas matemáticas sustantivas; el siguiente
conteo explica una de sus funciones sin reemplazarlas.

\begin{lema}[Rango y torsión en una imagen finita]
\label{lem:cohom-rango-torsion}
Sea $H$ un grupo abeliano cuyo subgrupo de torsión es
finito, de orden $T$. Sea $J\subset H$ de rango a lo
sumo $r$ y $f:H\to B$ un homomorfismo, donde $qB=0$
para un entero $q\ge1$. Entonces
$\#f(J)\le q^rT$.
\end{lema}

\begin{proof}
Elijamos un conjunto finito de valores distintos de
$f(J)$ y preimágenes en $J$. El subgrupo $J_0$ generado
por ellas es finitamente generado:
$J_0\simeq\Z^{r_0}\oplus T_0$, con $r_0\le r$ y
$|T_0|\le T$.
La restricción de $f$ factoriza por $J_0/qJ_0$, de
cardinal a lo sumo $q^rT$.
Por tanto, ningún conjunto finito de valores de $f(J)$
tiene más elementos. Esto prueba la finitud de la
imagen completa y la cota.
\end{proof}

En la aplicación, la reducción, el homomorfismo de
Bockstein y la estabilización aportan dos factores
de torsión de homología integral, además de un factor
$(\ell^\nu)^{r_2(L)}$.
El conteo de caracteres de clases contiene ese mismo
factor, que se cancela. El resultado de A, (8.7),
ya no depende de $\nu$.
Esta uniformidad es lo que permite pasar al límite
con una cota de longitud.
Saber que cada grupo de clases de cada capa es finito
no daría la misma conclusión.

La entrada de torsión utilizada debe incluir todos los
primos. A, p.~40, advierte que una cota clásica formulada
solo salvo primos pequeños no basta para ese propósito.
Cuando faltan ciertas condiciones de raíces de unidad
o de rango, se pasa a un cuerpo auxiliar con controles
específicos de grado, discriminante y descenso de la
longitud; no se presupone monotonía para cualquier
extensión de cuerpos.

Finalmente, la dualidad y el cálculo local dan, en la
situación de A, (8.8),
\begin{equation}
 W_\ell(L,e)
 \le s_{\ell,e}(L)
     +n\bigl(2+\lfloor\log_\ell n\rfloor+e\bigr),
 \qquad n=[L:\Q].
 \label{eq:cohom-correccion-local}
\end{equation}
La hipótesis de raíces implica $\ell\le n+1$.
La corrección local, multiplicada por $\log\ell$,
se absorbe aumentando el exponente de
\eqref{eq:cohom-lema-k}.
Los controles principales de tamaño proceden de las
entradas homológicas y de clases anteriores.
En esta entrega se coteja su ensamblaje y su aplicación;
no se reconstruye toda esa teoría.

\subsection{Dos cotas de torre y una base diádica uniforme}

Sea $L_\nu/L_0$ la capa de nivel $\nu$ de una torre
ciclotómica $\Z_\ell$ sobre una base abeliana apropiada.
La base contiene $\mu_\ell$, y $\mu_4$ en el caso diádico.
Las entradas de A, p.~40, distinguen dos estimaciones.

La primera proporciona un entero $R$ que acota el
número de generadores del módulo pertinente,
$R'\le R$ centros $s_1,\ldots,s_{R'}$ en una extensión
ultramétrica de $\Q_\ell$ y un entero $c$ que registra
raíces de unidad en la base. Ponemos
\[
 j_0=\lfloor\log_\ell R\rfloor+3\quad(R>0);
 \qquad j_0=3\quad(R=0).
\]
Esta última convención coincide con
\texttt{Nat.log} en el código.
Si $B\ge0$ es entero y
$|s_i-e|_\ell\ge\ell^{-B}$ para todo $i$, se tiene
\begin{equation}
 s_{\ell,e}(L_\nu)
 \le R'\bigl(\max\{\nu,j_0\}+c+B\bigr).
 \label{eq:cohom-torre-ordinaria}
\end{equation}
El valor absoluto en la base se normaliza por
$|\ell|_\ell=\ell^{-1}$.
La cota se aplica cuando el torcido está suficientemente
separado de los centros.

La segunda estimación no exige esa separación:
\begin{equation}
 s_{\ell,e}(L_\nu)
 \le W_\ell(L_{j_0},e)+R(\nu-j_0)_+,
 \qquad (u)_+=\max\{u,0\}.
 \label{eq:cohom-segunda-cota}
\end{equation}
Esta es la forma de \texttt{SecondBound} leída en Lean:
el primer término de la derecha es la longitud de
$H^2$ en una capa fija.
La prueba cotejada compara primero longitudes de Sha
entre capas y después usa
$s_{\ell,e}(L_{j_0})\le W_\ell(L_{j_0},e)$.
La resta del código es la resta natural truncada.

La desigualdad se formula inicialmente con longitudes
extendidas. La entrada \eqref{eq:cohom-lema-k} da
finitud y un número concreto para la capa fija; luego
\eqref{eq:cohom-segunda-cota} transporta ambos a la
capa variable.
Aplicar directamente la cota gruesa a toda la capa
variable introduciría su grado creciente elevado a
$C_K$ en el exponente.
El traslado a grado acotado es esencial para el
presupuesto final.
B, pp.~63--64, comprime este paso; A, pp.~40--41, y
el tipo formal permiten hacerlo explícito.

\begin{lema}[Modelo de crecimiento por capas]
\label{lem:cohom-crecimiento}
Sea $M$ un $\Z_\ell$-módulo generado por a lo sumo
$R$ elementos y $\lambda_\nu=\length(M/\ell^\nu M)$.
Entonces esos cocientes son finitos y, para $\nu,j\ge0$,
\[
 \lambda_\nu\le\lambda_j+R(\nu-j)_+.
\]
\end{lema}

\begin{proof}
El cociente de nivel $\nu$ es cociente de
$(\Z_\ell/\ell^\nu\Z_\ell)^R$, así que es finito.
Si $\nu\le j$, es también cociente del de nivel $j$.
Si $\nu\ge j$, la sucesión
\[
 0\longrightarrow\ell^jM/\ell^\nu M
 \longrightarrow M/\ell^\nu M
 \longrightarrow M/\ell^jM\longrightarrow0
\]
es exacta. El primer término está generado por a lo
sumo $R$ elementos y anulado por $\ell^{\nu-j}$;
es cociente de
$(\Z_\ell/\ell^{\nu-j}\Z_\ell)^R$ y tiene longitud
a lo sumo $R(\nu-j)$.
La aditividad termina la prueba.
\end{proof}

Este modelo demuestra cómo una cota de generadores
produce un resto lineal en los niveles añadidos.
La identificación de los objetos de Iwasawa con los
cocientes pertinentes y su comparación con Galois
son entradas adicionales.
Tampoco se identifica sin más la cota de generadores
con el rango de un módulo libre.

\paragraph{Uniformidad cuando varían los primos auxiliares.}
Para los primos impares de $\mathcal B_\psi$, las bases
usadas en la reducción se fijan por $\psi$ y $\ell$.
En el caso diádico, una base auxiliar depende de $q,d$.
La anulación clásica del invariante $\mu$ para un
cuerpo abeliano fijo, de Ferrero--Washington
\cite{ferrero-washington}, no proporciona por sí sola
una cota uniforme al variar esos dos primos.
La entrada necesaria tiene el orden
\[
 \text{fijado }\psi,\quad
 \exists R\quad\forall(q,d)\text{ admisibles}\quad
 \text{las estimaciones de torre se cumplen}.
\]
Los centros pueden depender de la pareja; su número
está acotado antes de elegirla.

La interfaz diádica \texttt{TwoBase} añade
$\mu_4\subset L_0$, $[L_0:\Q]\le8o$ y una inmersión
ciclotómica de conductor $4N_\psi qd$.
Para el cuerpo $E_a$ de \eqref{eq:cohom-campos}
se construye una capa auxiliar con
\begin{equation}
 \nu\le a+2,\qquad
 s_{2,e}(E_a)\le1+s_{2,e}(L_\nu).
 \label{eq:cohom-puente-diadico}
\end{equation}
Se mantiene el término $+1$: al sumar $1\le e<t$
cuesta a lo sumo $t\log2$.
La comparación no identifica literalmente $E_a$,
$L_0$ y $L_\nu$.
El control de raíces en la base proporciona también
una cota uniforme para $c$.

El cotejo de \texttt{TwoIn} conserva precisamente
la elección de $R$ antes de $q,d$.
En \texttt{towerTwo} se usa
$R=4(g+6\deg M)$, donde $g$ procede de una base fija,
y después se aplican pasos específicos de extensiones
de caracteres y extensiones cuadráticas.
La existencia de un número finito de generadores para
cada pareja, sin un control común, no bastaría.
Una vez fijados $R$ y $j_0$, el grado de $L_{j_0}$
está acotado; su discriminante todavía depende de $q,d$.

\subsection{Pocos torcidos excepcionales}

El paso siguiente es finito y no requiere teoría de
Iwasawa. Es el Lema~8.5 de A, p.~40.

\begin{lema}[Separación ultramétrica]
\label{lem:cohom-excepcionales}
Sean $s_1,\ldots,s_r$ elementos de una extensión
ultramétrica de $\Q_\ell$, con el valor absoluto usual
en la base, y $Z\ge1$ entero.
Existe $\mathcal E\subseteq\{1,\ldots,Z\}$ con
$|\mathcal E|\le r$ tal que
\[
 e\notin\mathcal E,\quad1\le e\le Z
 \quad\Longrightarrow\quad
 |s_i-e|_\ell\ge\ell^{-\lfloor\log_\ell Z\rfloor}
 \quad\text{para todo }i.
\]
\end{lema}

\begin{proof}
Para cada $i$ elegimos un entero $e_i$ del intervalo
que minimice $|s_i-e_i|_\ell$ y ponemos
$\mathcal E=\{e_1,\ldots,e_r\}$.
Si $e\ne e_i$, la desigualdad ultramétrica y la
minimalidad dan
\[
 |e-e_i|_\ell
 \le\max\{|s_i-e_i|_\ell,|s_i-e|_\ell\}
 =|s_i-e|_\ell.
\]
El entero no nulo $e-e_i$ tiene valor absoluto
ordinario a lo sumo $Z$, así que
$\ord_\ell(e-e_i)\le\lfloor\log_\ell Z\rfloor$.
Esto prueba la desigualdad, incluso si el mínimo es
cero. Si $r=0$, sirve el conjunto vacío.
\end{proof}

Se aplica una sola vez con $Z=\lfloor X^{401}\rfloor$.
El intervalo cubre todos los torcidos necesarios para
todos los $x\le X$.
Reunimos los conjuntos excepcionales de las torres
impares pertinentes y de la torre diádica.
Su unión tiene cardinal a lo sumo
\[
 R_\psi=R_2+\sum_{\ell\in\mathcal B_\psi}R_\ell,
\]
con las cotas de generadores de las torres usadas.
Los conjuntos y sus posiciones pueden depender de
$X,q,d$, pero $R_\psi$ ya está fijado.
Fuera de la unión se aplica
\eqref{eq:cohom-torre-ordinaria} con
$B=\lfloor\log_\ell Z\rfloor$.

La selección no controla todavía el tamaño ordinario
de los enteros excepcionales. El siguiente lema
elige una escala que proporciona exactamente ese
control para los que intervienen en las sumas.

\subsection{Una ventana común por un argumento de huecos}

\begin{lema}[Hueco multiplicativo]
\label{lem:cohom-ventana}
Sea $\mathcal E\subset\N_{>0}$ con $|\mathcal E|\le R$,
$0<\beta\le C_0$ y $X>1$.
Pongamos
\[
 J=2^{R+2}(C_0/\beta)^R.
\]
Existe $W$ con $X^{1/J}\le W$ y $W^2\le X$ tal que,
para todo $x\in[W,W^2]$ y $e\in\mathcal E$,
\[
 e\le x^{C_0}\quad\Longrightarrow\quad e\le x^\beta.
\]
\end{lema}

\begin{proof}
Escribamos $U=\log X$ y
\[
 u_i=\frac U2\left(\frac{\beta}{2C_0}\right)^i,
 \qquad 0\le i\le R.
\]
Los $R+1$ intervalos $(\beta u_i,2C_0u_i]$ son
disjuntos: el extremo superior del siguiente coincide
con el extremo inferior, excluido, del anterior.
Uno no contiene ninguno de los a lo sumo $R$ números
$\log e$, con $e\in\mathcal E$.
Elegimos ese índice y ponemos $W=\exp(u_i)$.
Entonces
\[
 \frac{\log X}{J}
 \le\frac{U}{2^{R+1}(C_0/\beta)^R}
 \le u_i\le U/2.
\]
Esto da las dos condiciones sobre $W$.
Si $W\le x\le W^2$ y
$x^\beta<e\le x^{C_0}$, se tendría
$\beta u_i<\log e\le2C_0u_i$, contradicción.
\end{proof}

Esta prueba sigue B, Lema~5.5, p.~63.
A, Lema~8.6, p.~41, prueba la selección equivalente
por inducción, suprimiendo un punto y reduciendo la
escala. El cotejo de \texttt{window1} siguió esa
inducción completa en Lean.
En ambos casos se conserva la misma constante suficiente
para la aplicación.

Tomamos ahora
\begin{equation}
 C_0=401,\qquad
 \beta=\frac{\varepsilon}{2C_K},\qquad
 J=2^{R_\psi+2}
       \left(\frac{802C_K}{\varepsilon}\right)^{R_\psi}.
 \label{eq:cohom-eleccion-j}
\end{equation}
La condición $C_K\ge1$ y el rango de $\varepsilon$
aseguran $\beta\le C_0$.
El resultado es la implicación simultánea
\begin{equation}
 e\in\mathcal E,\quad e\le x^{401},\quad W\le x\le W^2
 \quad\Longrightarrow\quad
 e\le x^{\varepsilon/(2C_K)}.
 \label{eq:cohom-torcido-pequeno}
\end{equation}
Los torcidos excepcionales no tienen que desaparecer.
Si están presentes a escala $x$, son suficientemente
pequeños para la cota de capa fija.
El intervalo se elige una sola vez para toda la unión
de excepciones, de modo que sirve simultáneamente
a todos los primos pertinentes.

\subsection{El discriminante y la cuenta final de exponentes}

Mostramos ahora cómo las entradas anteriores producen
\eqref{eq:cohom-meta}. Todas las constantes que siguen
dependen solo de los datos fijados y de $\varepsilon$;
los umbrales pueden aumentarse un número finito de veces
sin depender de la pareja admisible.

\paragraph{Conductor y discriminante.}
Puesto que $x\ge W\ge Y_0=X^{1/J}$,
\eqref{eq:cohom-primos} y \eqref{eq:cohom-altura} dan
\begin{equation}
 \log N_{a(x)}
 \le C_\psi\left[\log x+
              x^{\varepsilon/5}(1+\log x)^2\right].
 \label{eq:cohom-conductor}
\end{equation}
En las capas fijas pertinentes el grado está acotado
en términos de $\psi$.
El control de la inmersión ciclotómica y del
discriminante de la capa da una cota del mismo tipo:
\begin{equation}
 [L_{j_0}:\Q]\le D_\psi,\qquad
 \log(2|\operatorname{disc}L_{j_0}|)
 \le C_\psi\left[\log x+
              x^{\varepsilon/5}(1+\log x)^2\right].
 \label{eq:cohom-discriminante}
\end{equation}
La comparación de cuerpos y discriminantes es una entrada
de la torre, cotejada en \texttt{Layerdisc}.
El término que contiene $\log(qd)$ está conservado
mediante \eqref{eq:cohom-primos}.
Un grado acotado no implica que el discriminante
sea constante.

\paragraph{Los torcidos excepcionales.}
Por \eqref{eq:cohom-torcido-pequeno},
$(2e[L_{j_0}:\Q])^{C_K}\le C_\psi x^{\varepsilon/2}$.
El exponente del lado derecho de
\eqref{eq:cohom-lema-k} es, por tanto, a lo sumo
\begin{equation}
 \begin{split}
 C_\psi x^{\varepsilon/2}
       [\log x+x^{\varepsilon/5}(1+\log x)^2]
 \le C_\psi' x^{3\varepsilon/4}
 \end{split}
 \label{eq:cohom-exponente-final}
\end{equation}
para $x$ suficientemente grande.
En efecto,
$\varepsilon/2+\varepsilon/5=7\varepsilon/10
 <3\varepsilon/4$.
Al dividir por $x^{3\varepsilon/4}$ quedan
$x^{-\varepsilon/4}\log x$ y
$x^{-\varepsilon/20}(1+\log x)^2$, que tienden a cero.
La contribución de la capa fija es entonces
$\exp(C_\psi'x^{3\varepsilon/4})$ para la longitud
ponderada por $\log\ell$.
La comparación \eqref{eq:cohom-segunda-cota} añade
el resto lineal, y \eqref{eq:cohom-puente-diadico}
conserva el error uno.

\paragraph{Los torcidos ordinarios.}
En la ventana,
$\log X\le J\log x$ y
$\log Z\le401\log X$.
La altura diádica satisface
$\nu\le a(x)+2=O_\varepsilon(\log x)$.
Las otras torres requeridas por la reducción tienen
sus controles de nivel fijados por las entradas de
A, p.~40. Aplicando \eqref{eq:cohom-torre-ordinaria},
el coste por torcido ordinario en estas torres es
$O_{\psi,\varepsilon}(\log x)$.
Sumar hasta $t\le x^{401}$ y añadir el error diádico
de \eqref{eq:cohom-puente-diadico} cuesta a lo sumo
$C_{\psi,\varepsilon}x^{401}(1+\log x)$.

Hay a lo sumo $R_\psi$ torcidos excepcionales y un
número fijo de torres. Así, el total de las contribuciones
diádica y de torre excepcional está acotado por una
expresión de la forma
\[
 C_{\psi,\varepsilon}x^{401}(1+\log x)
 +C_{\psi,\varepsilon}
        \exp(C_\psi'x^{3\varepsilon/4}).
\]
El margen $3/4<4/5$ absorbe las constantes y el término
polinómico. Tras aumentar un umbral uniforme, esto
proporciona las cotas de B, p.~64,
\begin{equation}
 B_2\le\exp(x^{4\varepsilon/5}),\qquad
 B_{\rm exc}\le\exp(x^{4\varepsilon/5}).
 \label{eq:cohom-cotas-torres}
\end{equation}
El peso $\log2$ de $B_2$ se insertará al final;
pasar entre longitud y longitud ponderada introduce
solo constantes para estos primos fijos.

\paragraph{Cerrar la suma.}
La estimación \eqref{eq:cohom-conductor} y
$t=\lfloor x^{401}\rfloor$ implican, eventualmente,
\[
 \log(F_at)\le\log(N_at)\le2x^{\varepsilon/4}.
\]
El margen $\varepsilon/4-\varepsilon/5$ absorbe
los logaritmos y las constantes.
Como $n_E\le F_a$, la cota
\eqref{eq:cohom-menos-total} es a lo sumo
$C_1(F_at)^4$, con $C_1$ dependiente de $\psi$.
Para comprobar esta simplificación, $n_E^2t^2\le(F_at)^2$
y cada uno de los finitos factores
$\log(t\ell F_a)^2$ se domina por una constante por
$(F_at)^2$; el primer término se absorbe de igual modo.
Por tanto
\[
 B_-\le C_1\exp(8x^{\varepsilon/4}),\qquad
 B_+\le\exp(2Cx^{\varepsilon/4})
                  +C\exp(x^{4\varepsilon/5}).
\]
Insertando estas expresiones en
\eqref{eq:cohom-particion} obtenemos
\begin{equation}
 \begin{split}
 \mathcal V(E_{a(x)},t)
 &\le(\log2+C)\exp(x^{4\varepsilon/5})
       +C_1\exp(8x^{\varepsilon/4})
       +\exp(2Cx^{\varepsilon/4})\\
 &\le\exp(x^\varepsilon/2)
 \end{split}
 \label{eq:cohom-cierre}
\end{equation}
para un umbral uniforme suficientemente grande.
La última desigualdad se prueba dividiendo el logaritmo
de cada término por $x^\varepsilon$:
todos tienden a cero, y el número de términos es fijo.
Por ejemplo, se acota cada término por
$\exp(x^\varepsilon/4)$ y se absorbe su suma en
$\exp(x^\varepsilon/2)$ al aumentar el umbral.

Una potencia fija de $N_at$ no se ha declarado
polinómica en $x$; lo demostrado es
$\log((N_at)^C)=o(x^\varepsilon)$.
El objetivo es
$\exp(x^\varepsilon/2)=\sqrt{\exp(x^\varepsilon)}$.
La expresión $\exp(x^{\varepsilon/2})$ sería una meta
mucho menor y no se obtiene de esta comparación.

\begin{proposicion}[Ventana aritmética condicionada]
\label{prop:cohom-ventana-aritmetica}
Recibidas las entradas de congruencias, caracteres,
automorfía y compatibilidad local, la cota de capa fija
y los controles uniformes de torre descritos en esta
nota, se obtiene \eqref{eq:cohom-meta} con el orden de
elecciones allí indicado.
\end{proposicion}

\begin{proof}
Se fijan las cotas de generadores y $C_K$, luego $J$
por \eqref{eq:cohom-eleccion-j} y un umbral común
para las comparaciones anteriores.
Dado $X$ y una pareja admisible, el
Lema~\ref{lem:cohom-excepcionales} produce una unión
de excepciones de tamaño a lo sumo $R_\psi$.
El Lema~\ref{lem:cohom-ventana} elige una sola ventana.
La estimación \eqref{eq:cohom-cierre} vale para todos
sus puntos y para la misma pareja de primos.
Las entradas de torre y de caracteres proporcionan
finitud antes de convertir las longitudes en números;
los productos enteros y las cotas de sumas parciales
permiten aplicar el Lema~\ref{lem:cohom-presupuesto}.
Esto da simultáneamente la suma y su cota.
\end{proof}

Si el corte $h$ de \eqref{eq:cohom-entrada-indice}
satisface $h\le\lfloor x^{401}\rfloor$, la monotonía
del corte permite insertar esta estimación en el
defecto integral de los periodos.
La elección final de parámetros todavía debe comprobar
esa desigualdad y los demás rangos de la construcción.

\subsection{Concordancia localizada con Lean y alcance}

\begingroup
\raggedright
Se conserva el espejo
\texttt{CocoSgt/riemann-hypothesis-lean} en la revisión
\texttt{7cbf7a0c32fca1badf05931879a37c3f8a2b03b4}.
La concordancia de esta etapa comprende 36 archivos:
34 lecturas completas y dos selectivas.
Se recuperaron sus contenidos completos y se
recalcularon SHA-256 y hashes de blobs Git.
El manifiesto de \texttt{REGISTRO.json} registra las
rutas, los intervalos leídos, las declaraciones,
las pruebas y las hipótesis retenidas.
Un archivo homónimo descartado durante la búsqueda
queda fuera de ese recuento.

\begin{enumerate}
\item \textbf{Coeficientes, longitudes y dualidad.}
\href{https://github.com/CocoSgt/riemann-hypothesis-lean/blob/7cbf7a0c32fca1badf05931879a37c3f8a2b03b4/src/Formalism/GaloisCohomology/Statement.lean}{\texttt{GaloisCohomology/Statement}}
distingue \texttt{qzTw}, \texttt{Sha}, \texttt{shaLen}
y \texttt{VLen}.
Las longitudes permanecen en $\N\cup\{\infty\}$.
\href{https://github.com/CocoSgt/riemann-hypothesis-lean/blob/7cbf7a0c32fca1badf05931879a37c3f8a2b03b4/src/Development/AutomorphicArithmetic/AutomorphicVarieties/GaloisLibrary/GaloisCohomologyLength/Formula.lean}{\texttt{Formula}}
prueba la identidad bajo
\texttt{PoitouTateDuality}, \texttt{LocalFinitenessK},
\texttt{LocalTateDualityK} y el caso permitido
$\ell\ne2$ o ausencia de lugares reales.
\href{https://github.com/CocoSgt/riemann-hypothesis-lean/blob/7cbf7a0c32fca1badf05931879a37c3f8a2b03b4/src/Development/AutomorphicArithmetic/AutomorphicVarieties/AdditiveValuationNorms/ShaLeV.lean}{\texttt{ShaLeV}}
deduce la desigualdad de longitudes extendidas.
Las pruebas de dualidad recibidas por esos archivos
siguen siendo entradas.

\item \textbf{Las dos estimaciones de torre.}
\href{https://github.com/CocoSgt/riemann-hypothesis-lean/blob/7cbf7a0c32fca1badf05931879a37c3f8a2b03b4/src/Development/AutomorphicArithmetic/AutomorphicVarieties/LayeredIsomorphisms/SecondBound.lean}{\texttt{SecondBound}}
define proposiciones:
\texttt{ABound} reúne la cota fuera de centros y
la segunda cota.
En esta última, \texttt{VLen} corresponde a la capa
fija $j_0$, mientras \texttt{shaInv} corresponde a la
capa variable.
Se leyeron las pruebas de crecimiento y centros de
\texttt{Growth} y \texttt{ExistsCentresLevel},
y los ensamblajes
\href{https://github.com/CocoSgt/riemann-hypothesis-lean/blob/7cbf7a0c32fca1badf05931879a37c3f8a2b03b4/src/Development/AutomorphicArithmetic/AutomorphicVarieties/AdditiveValuationNorms/ABoundOfDescends.lean}{\texttt{ABoundOfDescends}}
y \texttt{ABoundOfIsTower}.
La teoría que suministra la torre, su descenso y
las comparaciones importadas no se auditó completa.
El Lema~\ref{lem:cohom-crecimiento} es un modelo
algebraico demostrado aquí; no se presenta como
transcripción de todo ese control.

\item \textbf{Cota de capa fija y uniformidad diádica.}
\href{https://github.com/CocoSgt/riemann-hypothesis-lean/blob/7cbf7a0c32fca1badf05931879a37c3f8a2b03b4/src/Development/AutomorphicArithmetic/AutomorphicVarieties/GeneratorCountTuples/LemmaK.lean}{\texttt{LemmaK}}
incluye la finitud de $H^2$ en su conclusión y presupone
el cuerpo totalmente complejo, $e\ge1$ y las raíces de
unidad de \eqref{eq:cohom-lema-k}.
Se siguió la ruta numérica de
\texttt{LemmaKOfSteps}, \texttt{NumThree} y
\texttt{SFinal}, y el paso de entradas importadas de
\texttt{LemmaKOfLitN}.
El cierre \texttt{lemmaK\_litN\_alt} de
\href{https://github.com/CocoSgt/riemann-hypothesis-lean/blob/7cbf7a0c32fca1badf05931879a37c3f8a2b03b4/src/Development/AutomorphicArithmetic/AutomorphicVarieties/ModularityArgument/ThresholdArguments/DenseThresholdCases/KFinal.lean}{\texttt{KFinal}}
conserva explícitamente
\texttt{LiteratureRankNHypothesis}.
En
\href{https://github.com/CocoSgt/riemann-hypothesis-lean/blob/7cbf7a0c32fca1badf05931879a37c3f8a2b03b4/src/Development/AutomorphicArithmetic/AutomorphicVarieties/QuadraticExtensions/IsAlgClosureRat.lean}{\texttt{QuadraticExtensions/IsAlgClosureRat}},
\texttt{towerTwo}, \texttt{twoBase} y \texttt{twoIn}
mantienen la elección uniforme de $R$, el control
de grado y el $+1$ diádico.
El cotejo distingue esas pruebas de los pasos de
Iwasawa y extensiones cuadráticas que importan.

\item \textbf{Excepciones, discriminante y una ventana.}
Se leyó completa la inducción de
\href{https://github.com/CocoSgt/riemann-hypothesis-lean/blob/7cbf7a0c32fca1badf05931879a37c3f8a2b03b4/src/Development/AutomorphicArithmetic/AutomorphicVarieties/LayeredIsomorphisms/Window1.lean}{\texttt{window1}},
con la constante $2^{R+2}B^R$.
\texttt{Exc} construye las excepciones mediante
vecindades ultramétricas; el
Lema~\ref{lem:cohom-excepcionales} sigue la versión
matemática de A con minimizadores.
\texttt{Layerdisc} conserva el discriminante de
$q,d$ y \texttt{Kexp} usa
$\varepsilon/2+\varepsilon/4=3\varepsilon/4$.
El teorema
\href{https://github.com/CocoSgt/riemann-hypothesis-lean/blob/7cbf7a0c32fca1badf05931879a37c3f8a2b03b4/src/Development/AutomorphicArithmetic/AutomorphicVarieties/LayeredIsomorphisms/ImportsAlt.lean}{\texttt{UpperTheory.V.Win.point}}
aplica \texttt{LemmaK} a la capa fija en la rama
excepcional y vuelve mediante \texttt{SecondBound}.
Prueba finitud antes de aplicar \texttt{toNat}.
\texttt{Point2}, \texttt{PartA}, \texttt{PartB} y
\texttt{ThmBabOf} mantienen los transportes, los
umbrales y una misma pareja auxiliar.
Se leyó ese ensamblaje condicionado, sin cerrar
todos los casos automorfos que lo alimentan.

\item \textbf{Sumabilidad antes del defecto integral.}
En
\href{https://github.com/CocoSgt/riemann-hypothesis-lean/blob/7cbf7a0c32fca1badf05931879a37c3f8a2b03b4/src/Formalism/SharpRestriction/GlobalTheoremInputs.lean}{\texttt{GlobalTheoremInputs}},
\texttt{ShaFinOdd} cuantifica todos los primos,
incluido dos, pese a su nombre.
\texttt{LinkN\_at} recibe expresamente
\texttt{Summable (vvTerm D)}.
Su suma de defectos es
$\sum_w\texttt{lwP}(w)\texttt{ramIdx}(w)
       \texttt{etaJ}(j,w)$, donde $w$ recorre el soporte
retenido con $\texttt{muAt}(w)\ne\texttt{mu0}$.
No incluye un factor dos adicional.
La cota tiene la forma de A, (7.9),
\[
 \texttt{etaSum}\le d_j^0
  \left[C\{fmbL^2+(\log x)^2N_c\}
                 +\sum_\ell'\texttt{vvTerm}(D,\ell)\right].
\]
Los símbolos no cohomológicos conservan los parámetros
de esa interfaz.
En
\href{https://github.com/CocoSgt/riemann-hypothesis-lean/blob/7cbf7a0c32fca1badf05931879a37c3f8a2b03b4/src/Development/AutomorphicArithmetic/CharacterBoxes/DivinNOfHypV.lean}{\texttt{DivinNOfHypV}},
\texttt{UpperTheory.Top.vvRecord} obtiene de
\texttt{VVle} la finitud individual, la sumabilidad
y la cota total.
Después \texttt{divinN\_of\_hypV} combina
\texttt{HypV}, \texttt{LinkN} y
\texttt{SmallLevelRankEq}.
La igualdad de rangos es una entrada separada.
No se ha reconstruido en esta lectura toda la
identificación de \texttt{etaJ} con la longitud
normalizada por ramificación de B.
\end{enumerate}

Las lecturas selectivas fueron
\texttt{GlobalTheoremInputs.lean}, líneas 1--63, y
\texttt{ExportBoundOfAdmissible.lean}, líneas 1--83
y 1000--1164. Se examinó el tramo final de exportación,
sin declarar revisados \texttt{point\_bound} ni el
desarrollo completo de primos grandes y excepcionales.
Las construcciones integrales de todos los proyectores,
la clausura de homología, torsión, dualidad y automorfía,
y la identidad íntegra del espejo con la distribución
original de Alpöge quedan fuera del alcance.
No se ejecutó Lean ni se obtuvo un informe de axiomas.

El avance matemático de la nota consiste en las pruebas
finitas y en el ensamblaje de la ventana bajo entradas
explícitas. Permite seguir el coste cohomológico hasta
el defecto integral de los periodos.
El siguiente desarrollo reunirá los presupuestos en
una elección uniforme de parámetros, conservará las
entradas automorfas y de compatibilidad de A, §9,
y estudiará el enlace con las ventanas de $H_\psi$
en A, §§10--11, y B, §6.
La construcción aritmética completa y la propuesta
global no quedan certificadas por este cotejo.
\par
\endgroup

\subsection{Ejercicios con soluciones}

\begin{ejercicio}\label{ej:cohom-congruencias}
Sea $\mathcal O=\Z_\ell$, $r,s\ge1$ y
\[
 M=\mathbb T
 =\{(x,y,z)\in\mathcal O^3:
          x\equiv y\pmod{\ell^r},\
          y\equiv z\pmod{\ell^s}\},
\]
con acción por multiplicación coordenada.
Retenga las dos primeras coordenadas.
Calcule $e_ZM$, $M\cap V_Z$, los ideales
$\mathfrak c_1,\mathfrak c_2$ y la longitud del
cociente. Compare con la
Proposición~\ref{prop:cohom-congruencias}.
\end{ejercicio}

\paragraph{Solución del Ejercicio~\ref{ej:cohom-congruencias}.}
Omitiendo la tercera coordenada, que es cero,
\[
 L=\{(x,y):x-y\in\ell^r\mathcal O\},\qquad
 N=\{(x,y)\in L:y\in\ell^s\mathcal O\}.
\]
En $L$ se puede elegir $z=y$ para levantar cualquier
par a $M$. En $N$ se impone $z=0$.
La aplicación $(x,y)\mapsto y\bmod\ell^s$ induce
$L/N\simeq\mathcal O/\ell^s\mathcal O$, así que la
longitud es $s$.

Los elementos de $\mathfrak J$ son las ternas con $z=0$.
Su segunda coordenada recorre $\ell^s\mathcal O$;
la primera recorre
$\ell^s\mathcal O+\ell^r\mathcal O
 =\ell^{\min(r,s)}\mathcal O$.
Por tanto
\[
 \mathfrak c_1=\ell^{\min(r,s)}\mathcal O,\qquad
 \mathfrak c_2=\ell^s\mathcal O.
\]
La cota de la proposición es $s\le\min(r,s)+s$, estricta
cuando $r,s\ge1$.
El álgebra separa las coordenadas, por ejemplo mediante
múltiplos suficientemente divisibles de los vectores
coordenados. El ejemplo muestra que la estimación por
ideales no tiene por qué ser una igualdad.

\begin{ejercicio}\label{ej:cohom-normalizacion}
Tome $A=\Z_\ell/\ell^2\Z_\ell\oplus
\Z_\ell/\ell^5\Z_\ell$ y un anillo de enteros
$\mathcal O$ con ramificación tres y grado residual dos.
Calcule la longitud de $A$, la de $A\otimes\mathcal O$,
su longitud normalizada y el logaritmo de su cardinal.
Verifique también la exactitud y la cuenta de longitudes de
\[
 0\longrightarrow\Z/\ell\Z
 \xrightarrow{\ \ell^2\ }\Z/\ell^3\Z
 \xrightarrow{\ \ell^2\ }\Z/\ell^4\Z
 \longrightarrow\Z/\ell^2\Z\longrightarrow0,
\]
donde el último mapa es reducción.
\end{ejercicio}

\paragraph{Solución del Ejercicio~\ref{ej:cohom-normalizacion}.}
La longitud inicial es $2+5=7$.
Tras extender coeficientes se obtiene
$\mathcal O/\varpi^6\mathcal O\oplus
\mathcal O/\varpi^{15}\mathcal O$,
de longitud $21$.
Su longitud normalizada es $21/3=7$ y su logaritmo
de cardinal es $2\cdot21\log\ell=42\log\ell$.
La ramificación y el grado residual desempeñan
funciones distintas.

En la sucesión, el núcleo del segundo mapa son los
múltiplos de $\ell^2$ módulo $\ell^3$, imagen del
primero. Su imagen son los múltiplos de $\ell^2$
módulo $\ell^4$, núcleo de la reducción final.
Los extremos son inyectivo y sobreyectivo,
respectivamente, y los mapas están bien definidos.
La aditividad da $4+1=2+3$.
Es un modelo finito de la corrección de dualidad:
los términos análogos a $H^2$ y Sha tendrían longitudes
cuatro y dos, no iguales.
No se afirma que estos grupos concretos provengan
de un cuerpo de números.

\begin{ejercicio}\label{ej:cohom-crecimiento}
Para $M=\Z_\ell^2\oplus\Z_\ell/\ell^5\Z_\ell$,
calcule $\lambda_\nu=\length(M/\ell^\nu M)$ y
compruebe el Lema~\ref{lem:cohom-crecimiento} con
$R=3$, $j=2$ y $\nu=8$.
Explique por qué la familia $M_N=\Z_\ell^N$ tiene
cada módulo finitamente generado pero no una cota
uniforme del número de generadores.
\end{ejercicio}

\paragraph{Solución del Ejercicio~\ref{ej:cohom-crecimiento}.}
Los dos sumandos libres aportan $2\nu$ y el de torsión
aporta $\min(\nu,5)$. Por tanto
\[
 \lambda_\nu=2\nu+\min(\nu,5),\qquad
 \lambda_2=6,\qquad\lambda_8=21.
\]
La estimación es $21\le6+3(8-2)=24$.
Después del nivel cinco, el sumando de torsión deja
de crecer, por lo que una cota con tres generadores
no tiene que ser exacta.
Por otro lado, $M_N/\ell M_N$ tiene dimensión $N$
sobre $\mathbf F_\ell$, de modo que $M_N$ necesita
al menos $N$ generadores.
La afirmación $\forall N\,\exists R_N$ no implica
$\exists R\,\forall N$.
Es el orden de cuantificadores que debe verificarse
en la entrada diádica; el ejemplo no sustituye su
construcción aritmética.

\begin{ejercicio}\label{ej:cohom-ventana}
Aplique la prueba del Lema~\ref{lem:cohom-ventana}
a $\mathcal E=\{2,2^{16}\}$, $X=2^{16}$,
$C_0=4$, $\beta=1$ y $R=2$.
Encuentre una ventana y compruebe la implicación
para sus dos elementos, incluyendo los extremos.
\end{ejercicio}

\paragraph{Solución del Ejercicio~\ref{ej:cohom-ventana}.}
Se tiene $U=16\log2$, y los tres valores son
$u_0=8\log2$, $u_1=\log2$ y
$u_2=(\log2)/8$.
El primer intervalo $(8,64]\log2$ contiene
$\log(2^{16})$.
El segundo, $(1,8]\log2$, no contiene ninguno de los
dos logaritmos: $\log2$ es su extremo inferior
excluido. Elegimos $W=e^{u_1}=2$.

En $2\le x\le4$, el excepcional $2$ está presente
y satisface $2\le x=x^\beta$.
En cambio, $2^{16}>4^4=256$, por lo que nunca
satisface $e\le x^{C_0}$ en esa ventana.
No es necesario eliminar el primer excepcional.
La constante del lema es
$J=2^{2+2}4^2=256$:
$X^{1/J}=2^{1/16}\le2$ y $W^2=4\le X$.
Los extremos semiabiertos de los intervalos
logarítmicos son los que permiten esta elección.

\subsection{Lectura complementaria: conúcleos torcidos en una torre}
\label{subsec:cohom-fuente-c}

\fuente{C, \emph{Twisted cokernels in a cyclotomic tower},
«Deformalization 1», ejemplar de 22 páginas:
§§2--6, pp.~2--9, y §§7--8, pp.~9--15
\cite{fuenteC}. Esta lectura complementaria desarrolla el mecanismo
algebraico de las dos cotas de torre utilizadas en esta nota.}

La fuente C aporta una exposición localizada del paso de un módulo
de Iwasawa a las cotas de longitudes torcidas. Conservamos nuestra
notación: $e$ es el torcido y $\nu$ la capa; corresponden a $h$ y $m$
en C. El nivel del escalar se escribirá $a_u$, para distinguirlo
de la altura diádica $a$ utilizada antes.

\paragraph{Las hipótesis y los cuatro argumentos.}
Fijemos un primo $\ell$, un entero $R\ge0$, una sobreyección
$\Z_\ell^R\twoheadrightarrow X$
y $\gamma\in\operatorname{End}_{\Z_\ell}(X)$ tal que
$(\gamma-1)^N X\subseteq\ell X$ para algún $N\ge1$.
Sea $u\in\Z_\ell$ con
$a_u=v_\ell(u-1)$ entero finito, $a_u\ge1$, y $a_u\ge2$ si $\ell=2$.
No suponemos que $X$ sea libre: $R$ cuenta generadores, incluida
la posible torsión. Normalizamos $v_\ell(\ell)=1$, ponemos
$v_\ell(0)=\infty$ y
\[
 \tau(\gamma,u;e,\nu)=
 \length_{\Z_\ell}\operatorname{coker}
       \bigl((u^e\gamma)^{\ell^\nu}-1\bigr),
 \qquad e,\nu\in\N.
\]
La función está definida sin las hipótesis de unipotencia y nivel;
estas se necesitan para las cotas. Sus valores pertenecen a
$\N\cup\{\infty\}$. Como antes,
$j_0=\lfloor\log_\ell R\rfloor+3$ para $R>0$ y $j_0=3$ para $R=0$.
En este último caso $X=0$ y todas las longitudes son cero.

\paragraph{Centros comunes y crecimiento por capas.}
Los Teoremas A y B de C, p.~3, afirman que existen
$c_1,\ldots,c_R\in\Z_\ell$ tales que, para todos $e,\nu\ge0$,
\[
 \begin{aligned}
 \tau(\gamma,u;e,\nu)
 &\le\sum_{i=1}^R
       \bigl(\max\{\nu,j_0\}+a_u+v_\ell(c_i-e)\bigr),\\
 \tau(\gamma,u;e,\nu)
 &\le\tau(\gamma,u;e,j_0)+R(\nu-j_0)_+.
 \end{aligned}
\]
Primero se fijan $X,R,\gamma,u$ y después se eligen los centros,
que sirven para \emph{todos} los torcidos y todas las capas.
Si se cambia $\gamma$ o $u$, la existencia de centros para cada
nuevo par no proporciona centros comunes entre esos pares.

La construcción explica esta independencia. A una capa fija se
levanta $\gamma^{\ell^{j_0}}$ a una matriz sobre $\Z_\ell^R$.
El conúcleo matricial se sobreyecta al conúcleo de $X$, por lo que
su valoración del determinante da una \emph{cota superior}.
La identidad
$v_\ell(u^n-1)=a_u+v_\ell(n)$ y la compacidad de $\Z_\ell$
permiten asociar un centro a cada autovalor del levantamiento.
La desigualdad ultramétrica controla la proximidad a ese centro;
no se utiliza un logaritmo $\ell$-ádico ni se requiere que $X$
sea libre (C, §§3--4, pp.~5--7).

Para pasar a otras capas se usa un argumento distinto.
Si $w\equiv1\pmod{\ell^2}$ en el sentido de
$(w-1)X\subseteq\ell^2X$, entonces
$1+w+\cdots+w^{\ell-1}\equiv\ell\pmod{\ell^2}$.
Nakayama muestra que su imagen contiene $\ell X$; su conúcleo
tiene longitud a lo sumo $R$.
La factorización de $w^\ell-1$ limita cada aumento de capa a $R$.
La unipotencia módulo $\ell$ permite llegar a esa situación
en una capa acotada por $\ell$ y $R$.
Se transportan así los centros elegidos en $j_0$ a todos los niveles
(C, §5, pp.~7--9).

Para $1\le e\le Z$, cada centro admite a lo sumo un entero
con $v_\ell(c_i-e)>\lfloor\log_\ell Z\rfloor$.
Es el mecanismo del Lema~\ref{lem:cohom-excepcionales};
el Corolario C de la nueva fuente, pp.~3 y 9, obtiene
a lo sumo $R$ excepciones.
Su número es uniforme en $Z$, aunque sus posiciones puedan cambiar.

\paragraph{El puente aritmético y la capa fija.}
Sea $L_\infty/L$ la torre ciclotómica, con $[L_\nu:L]=\ell^\nu$
y con $L$ conteniendo
$\mu_\ell$, y $\mu_4$ si $\ell=2$.
El módulo $X_\infty$ considerado por C corresponde a la extensión
abeliana pro-$\ell$ máxima de $L_\infty$ no ramificada en ningún
lugar y en la que los primos sobre $\ell$ se descomponen completamente.
Supongamos que está generado por $R\ge1$ elementos sobre $\Z_\ell$.
Si $\chi_{\mathrm{cyc}}$ es el carácter ciclotómico, elegimos un
generador topológico $\sigma$ de $\operatorname{Gal}(L_\infty/L)$.
Entonces $u=\chi_{\mathrm{cyc}}(\sigma)$ genera su imagen;
$\gamma$ es la acción de $\sigma$ en $X_\infty$.
El descenso proporciona, para $e\ge1$,
\[
 s_{\ell,e}(L_\nu)=\tau(\gamma,u;e,\nu).
\]
La segunda cota y Poitou--Tate dan entonces
\[
 s_{\ell,e}(L_\nu)
 \le
 \underbrace{\length_{\Z_\ell}
 H^2(G_{L_{j_0},S_\ell},\Z_\ell(e+1))}_{W_\ell(L_{j_0},e)}
 +R(\nu-j_0)_+.
\]
Es precisamente la estructura de \eqref{eq:cohom-segunda-cota}:
el primer término pertenece a la \emph{capa fija}.
Una cota por centros puede ser infinita incluso si la longitud
efectiva es finita, pues el levantamiento matricial puede introducir
esa pérdida cuando hay torsión.
La finitud en todas las capas se obtiene de esta segunda desigualdad
y de la entrada finita para $H^2$, no del argumento de centros.

C distingue tres entradas en §7, pp.~10--11:
nilpotencia local sobre el módulo dual discreto, descenso de Sha
y comparación por Poitou--Tate.
La primera y la tercera se enuncian con bosquejos.
El descenso se desarrolla en §8, pero recibe los Inputs 2a y 2b
de p.~15: la identificación cohomológica con invariantes y su
reducción al núcleo de un generador.
La condición $e\ne0$ entra en la primera identificación;
el uso aquí requiere $e\ge1$.
Las cotas cuantitativas de la capa fija y la uniformidad de $R$
en familias conservan el carácter de entradas de esta nota.

\begin{center}
\small
\begin{tabularx}{\textwidth}{@{}p{0.30\textwidth}Y@{}}
\toprule
Pasaje de C & Lugar de estas notas \\
\midrule
Teoremas A--B; §§3--6 &
Las dos cotas de torre y el
Lema~\ref{lem:cohom-excepcionales}.\\
Teorema D; §§7--8 &
El núcleo local, la dualidad y
\eqref{eq:cohom-segunda-cota}.\\
§7, pp.~11--12 &
La entrada de capa fija \eqref{eq:cohom-lema-k},
la ventana y el orden de elecciones de la
Nota~\ref{sec:parametros}.\\
\bottomrule
\end{tabularx}
\end{center}

\paragraph{Procedencia y alcance del cotejo.}
C remite al espejo de Deicyde, revisión \texttt{40e56bab};
nuestras concordancias previas fijan el espejo de CocoSgt,
revisión \texttt{7cbf7a0}.
Las comprobaciones declaradas por C se atribuyen a la fuente;
esta incorporación no ejecuta Lean.
C, §10, p.~17, aclara que su formulación de los centros
en $\Z_\ell$, con el nivel exacto $a_u$, concreta más que el tipo
literal \texttt{ABound} usado antes.
Conservamos $j_0=\lfloor\log_\ell R\rfloor+3$:
la mejora con $+2$ que C observa no es la formulación del archivo.
La Nota~\ref{sec:finitud-uniforme} justifica la finitud para datos fijos;
\texttt{LemmaK}, \texttt{PairEx} y el cierre global conservan sus entradas.
